\documentclass{article}
\usepackage{graphicx}
\usepackage{amsmath}
\usepackage{booktabs}
\usepackage[margin=1in]{geometry}

\title{CSCI 311.S26 -- DNA Sequence Matching System}
\author{Anna Barrios}
\date{March 2026}

\begin{document}

\maketitle

\section{Introduction}

This project develops a DNA sequence similarity search tool. Given a query sequence $t$ and a collection of known sequences $D$ stored in FASTA format, the program compares $t$ against each sequence in $D$ using a selected algorithm and returns the most similar result.

The implementation is written in Python 3 and is organized into three main components: a file parsing module for handling FASTA input, a set of dynamic programming algorithms for sequence comparison, and a command-line interface that allows the user to select an algorithm and run the search.

The overall approach follows a simple structure: for every sequence $s_i$ in the database, compute a similarity score with $t$, track the highest score, and return the corresponding sequence. Each algorithm operates in $O(mn)$ time for sequences of lengths $m$ and $n$.

\section{Algorithms and Analysis}

\subsection{Longest Common Substring}

This method identifies the longest contiguous segment that appears in both sequences. A dynamic programming table is used where each entry stores the length of the matching substring ending at specific indices.

\[
dp[i][j] =
\begin{cases}
dp[i-1][j-1] + 1 & \text{if } s[i-1] = t[j-1] \\
0 & \text{otherwise}
\end{cases}
\]

The maximum value in the table represents the result.

\textbf{Time:} $O(mn)$ \quad \textbf{Space:} $O(\min(m,n))$

\textbf{Strength:} Efficient and good at detecting strong local matches.

\textbf{Limitation:} Requires exact contiguous matches. Even a single mismatch breaks the sequence, which is unrealistic for biological data.

\subsection{Longest Common Subsequence (LCS)}

The LCS algorithm allows matches that are not contiguous, as long as they appear in the same order. This is computed using dynamic programming with the recurrence:

\[
dp[i][j] =
\begin{cases}
dp[i-1][j-1] + 1 & \text{if } s[i-1] = t[j-1] \\
\max(dp[i-1][j], dp[i][j-1]) & \text{otherwise}
\end{cases}
\]

\textbf{Time:} $O(mn)$ \quad \textbf{Space:} $O(\min(m,n))$

\textbf{Strength:} More flexible than substring since it tolerates gaps.

\textbf{Limitation:} Does not penalize gaps, which can lead to inflated similarity scores for longer sequences.

\subsection{Edit Distance}

Edit distance measures how many operations (insertions, deletions, substitutions) are needed to transform one sequence into another. The recurrence is:

\[
dp[i][j] =
\begin{cases}
dp[i-1][j-1] & \text{if } s[i-1] = t[j-1] \\
1 + \min(dp[i-1][j], dp[i][j-1], dp[i-1][j-1]) & \text{otherwise}
\end{cases}
\]

To align with the system (higher = more similar), the value is negated.

\textbf{Time:} $O(mn)$ \quad \textbf{Space:} $O(\min(m,n))$

\textbf{Strength:} Captures all types of mutations and handles sequences of different lengths well.

\textbf{Limitation:} Assigns equal cost to all operations, which does not reflect biological differences between mutations.

\subsection{Global Alignment (Needleman-Wunsch)}

This algorithm computes an optimal global alignment using a scoring system. We use:
\begin{itemize}
    \item Match: +1
    \item Mismatch: -1
    \item Gap: -2
\end{itemize}

The recurrence is:

\[
dp[i][j] = \max
\begin{cases}
dp[i-1][j-1] + \text{score} & \text{(match/mismatch)} \\
dp[i-1][j] - 2 & \text{(gap in } t\text{)} \\
dp[i][j-1] - 2 & \text{(gap in } s\text{)}
\end{cases}
\]

\textbf{Time:} $O(mn)$ \quad \textbf{Space:} $O(mn)$

\textbf{Strength:} Produces the most realistic alignment by distinguishing between matches, mismatches, and gaps.

\textbf{Limitation:} Uses more memory and is slower compared to other approaches.

\section{System Structure}

The system processes each sequence independently using the selected algorithm. For each sequence in the dataset, a similarity score is computed and compared to the current best score. The highest-scoring sequence is returned at the end of the search.

This follows the standard formulation:
\[
\max_{s_i \in D} \text{Similarity}(s_i, t)
\]

\section{Discussion}

All implemented methods have the same asymptotic time complexity, but they differ significantly in how they evaluate similarity.

Simpler methods such as substring and LCS are faster and easier to implement but fail to fully capture biological relationships. Edit distance improves accuracy by modeling sequence transformations, while global alignment provides the most meaningful comparison by incorporating penalties for gaps and mismatches.

Overall, the project highlights the tradeoff between computational efficiency and biological accuracy. More sophisticated algorithms produce better results but require greater computational resources.

\end{document}
